% Part: computability
% Chapter: recursive-functions
% Section: primes

\documentclass[../../../include/open-logic-section]{subfiles}

\begin{document}

\olfileid[fr]{cmp}{rec}{pri}
\olsection{Nombres premiers}

La quantification bornée et la minimisation bornée nous donnent
de nombreux moyens de montrer que des fonctions et des relations
usuelles sont récursives primitives. Considérons, par exemple,
la relation « $x$ divise $y$ », notée $x \mid y$. Elle vaut lorsque
$y$ est un multiple entier de~$x$; dans le cas contraire, nous
écrivons $x \nmid y$. Autrement dit, $x \mid y$ si et seulement
s'il existe un~$z$ tel que $x\cdot z=y$. Un tel~$z$, s'il existe,
peut être choisi inférieur ou égal à~$y$. Il suffit donc de
chercher parmi les $z\leq y$. À partir de l'égalité et de la
multiplication, nous pouvons définir la divisibilité par une
quantification existentielle bornée :
\[
x \mid y \defiff \bexists{z \leq y}{(x \cdot z) = y}.
\]
Nous avons ainsi montré que la relation $x\mid y$ est récursive
primitive.\footnote{Corrections éditoriales : la phrase source sur le reste inverse le dividende et le diviseur ; nous définissons plutôt $x\nmid y$ comme la négation de $x\mid y$, y compris si $x=0$. La borne sur le témoin signifie qu'on peut en choisir un avec $z\leq y$, et non que tout témoin est borné : pour $x=y=0$, tout $z$ convient.}

Un entier naturel~$x$ est \emph{premier} s'il n'est ni $0$ ni
$1$ et si ses seuls diviseurs sont $1$ et lui-même. Autrement
dit, si $y\mid x$, alors $y=1$ ou $y=x$. Pour vérifier que~$x$
est premier, il suffit d'examiner les $y\leq x$ : lorsque
$x\geq2$, aucun $y>x$ ne divise~$x$. La relation
$\fn{Prime}(x)$, qui vaut exactement lorsque $x$ est premier,
se définit donc par
\[
\fn{Prime}(x) \defiff x \geq 2 \land \bforall{y \leq x}{(y \mid x \lif y
  = 1 \lor y = x)}
\]
et elle est récursive primitive.

Les nombres premiers sont $2$, $3$, $5$, $7$, $11$, etc.
Considérons la fonction $p(x)$ qui renvoie le premier de rang~$x$
dans cette suite : $p(0)=2$, $p(1)=3$, $p(2)=5$, etc. Nous
écrirons souvent $p_x$ pour $p(x)$; ainsi $p_0=2$, $p_1=3$,
et ainsi de suite.

Si nous disposions d'une fonction $\fn{nextPrime(x)}$ qui
renvoie le premier nombre premier strictement supérieur
à~$x$, nous pourrions définir~$p$ par récurrence primitive :
\begin{align*}
  p(0) & = 2\\
  p(x+1) & = \fn{nextPrime}(p(x))
\end{align*}
La fonction $\fn{nextPrime}(x)$ renvoie le plus petit~$y$
tel que $y>x$ et que $y$ soit premier. Une recherche non
bornée permettrait de la calculer, mais le théorème d'Euclide
permet de recourir à la minimisation bornée : il existe toujours
un nombre premier strictement supérieur à~$x$ et inférieur
ou égal à $\fact{x}+1$. Ainsi,
\[
  \fn{nextPrime(x)} =
  \bmin{y \leq \fact{x}+1}{(y > x \land \fn{Prime}(y))}.
\]
La borne est ici inclusive; on peut la remplacer par la borne
stricte $y<\fact{x}+2$ dans le schéma précédent. Il s'ensuit
que $\fn{nextPrime}(x)$, puis $p(x)$, sont récursives
primitives, et non simplement calculables.

(Voici une preuve rapide du fait utilisé. Pour $x=0$ ou
$x=1$, le nombre $2$ convient. Supposons désormais $x\geq2$.
Soit $p_n$ le plus grand nombre premier inférieur ou égal
à~$x$, et soit $p=p_0\cdot p_1\cdot\dots\cdot p_n$ le produit
de tous les nombres premiers inférieurs ou égaux à~$x$.
Aucun des nombres premiers inférieurs ou égaux à~$x$ ne divise
$p+1$ : chacun divise~$p$, tandis que la division de $p+1$ par ce
nombre laisse le reste~$1$.
Si $p+1$ est premier, il est donc strictement supérieur à~$x$.
Sinon, il possède un diviseur premier~$q<p+1$; d'après ce qui
précède, $q>x$. Enfin, $p\leq\fact{x}$; dans les deux cas, nous
avons trouvé un nombre premier strictement supérieur à~$x$
et inférieur ou égal à $\fact{x}+1$.)\footnote{Complément éditorial : la source commence par le plus grand premier inférieur ou égal à $x$. Celui-ci n'existe pas pour $x=0,1$ ; nous traitons ces deux cas séparément par le témoin $2$.}

\begin{prob}
Définir la division entière $d(x,y)$ au moyen de la
minimisation bornée.
\end{prob}

\end{document}
